% Lecture visuelle directe avec la vision du modèle ; aucun OCR.
% Source : PDF original, vue 005, côté droit, page manuscrite 9.
% Statut : transcrit ; relecture indépendante à effectuer.
\part*{Précis historique sur les causes qui ont amené la fondation de la Congrégation des jeunes gens en 1802}
\addcontentsline{toc}{part}{Précis historique sur les causes qui ont amené la fondation de la Congrégation des jeunes gens en 1802}
\chapter{Coup d'œil sur l'état de la France et de la ville de Lyon en particulier, depuis le commencement de la Révolution jusqu'au 9 Thermidor}

Je suis forcé, Messieurs et amis, de prendre les choses de bien haut, car pour pouvoir apprécier dans toute leur étendue les grâces que le Seigneur a répandues avec tant d'abondance sur les Congréganistes de Lyon, il est indispensable de jeter au moins un coup d'œil rapide sur la situation où l'impiété et les résultats de ses infernales manœuvres avaient placé en France et particulièrement dans notre ville la religion et la société.
% Le mot « triste » devant « situation » est barré dans la source.

Vous n'ignorez pas plus que moi, Messieurs, que le démon de l'orgueil a proclamé, il y a plus de trois siècles, par la bouche de ses suppôts, Luther, Calvin et consorts, l'indépendance de la pensée en matière de religion. Voulant étendre encore son œuvre de licence, le même esprit de mensonge s'est servi à cette fin, dans le 18\up{e} siècle, de ceux qui se sont eux-mêmes honorés du nom de philosophes.

\section*{Coup d'œil sur la Révolution}
% Mention marginale à partir de la page manuscrite 10 ; placée au début du paragraphe.
C'en a été trop peu d'attaquer seulement quelques vérités. Fidèles aux suggestions de cet ancien serpent qui leur a persuadé comme autrefois à Ève qu'avec le fruit qu'il leur présentait toute science leur serait acquise,
% Source : PDF original, vue 006, côté gauche, page manuscrite 10.
et qu'ils seraient comme des dieux, ces nouveaux précepteurs du genre humain ont voulu qu'à l'avenir, pour peser les vérités de la foi, le poids du sanctuaire fût remplacé par celui de leur intelligence. Ils ont livré la parole immuable sortie de la bouche même du Très-Haut à tout l'arbitraire des systèmes que leurs imaginations déréglées ont été capables d'inventer. C'est ainsi que quelques-uns d'entre eux ont bien daigné faire la grâce à Dieu de l'admettre dans le symbole de leur croyance, mais à condition toutefois que ce serait un dieu commode qui se bornerait à régner dans le ciel, sans se mêler en aucune manière des affaires de ce monde ; que surtout il cesserait de les effrayer par ces peines éternelles de l'enfer, dont ils ne veulent absolument pas, car ils les redoutent trop pour y croire. D'autres, plus hardis, n'ont plus voulu reconnaître d'autre dieu que la matière. Quelques-uns poussant les conséquences logiques de leurs principes jusqu'à la dernière extrémité ont fini par rejeter entièrement ces notions du juste et de l'injuste que Dieu par la loi naturelle a gravées dans le cœur de tous les hommes en caractères ineffaçables.

Le même esprit de licence qui portait ces novateurs à secouer le joug de la religion devait aussi les amener à se soustraire à l'autorité des lois. Un grand complot fut donc forgé de renverser en même temps, si cela eût été au pouvoir de ces hardis conspirateurs, et le Dieu qui règne dans le ciel, et tous ceux qui le représentent sur la terre.

Ici, Messieurs, permettez-moi une réflexion. Portez, je vous prie, avec moi les yeux plus haut. \textit{Sursum corda.}\footnote{Traduction : « Élevons notre cœur. » Référence : liturgie de la messe, dialogue de la préface. (Note du transcripteur).} Ne nous
% Source : PDF original, vue 006, côté droit, page manuscrite 11.
arrêtons pas seulement aux causes secondes ; mais considérons plutôt comment la Providence, dont la sagesse infinie dirige tous les événements, permet dans sa justice que le châtiment si justement mérité par l'orgueil insolent de l'impie se rencontre précisément dans le succès de ses projets liberticides. Un roi philosophe l'avait dit : Si je voulais punir une de mes provinces, je la donnerais à gouverner aux philosophes ; eh bien, pour punir la France qui a applaudi et exalté ce Voltaire, chef d'une philosophie anti-chrétienne et ses nombreux complices, Dieu dans sa colère va la livrer au gouvernement des disciples formés à cette école d'impiété. Mais allons les voir à l'œuvre, et nous pourrons alors juger de la nature de l'arbre par les fruits amers qu'il aura produits.

Ces artisans de ruine marchent tête levée au but qu'ils se sont proposé. Déjà ils ont renversé l'ordre des Jésuites. Il fallait bien commencer par détruire ce boulevard redoutable qui pouvait garantir la religion des coups qu'ils se préparaient à lui porter ; déjà ils ont travaillé, et ils n'ont que trop bien réussi à corrompre les mœurs, car ils n'ignorent pas que le désordre dans la conduite amène le désordre dans les idées, et que presque toujours l'impiété est fille du libertinage. Tout dans ce moment fait pressentir une catastrophe imminente. Les ministres du Dieu vivant, le cœur oppressé de douleur, la prédisent du haut de la chaire de vérité, et dans les transports d'une joie satanique les suppôts de l'enfer la nomment au milieu de leurs orgies.

Un seul obstacle pouvait retenir les efforts des méchants, mais
% Source : PDF original, vue 007, côté gauche, page manuscrite 12.
combien il sera facile à renverser. Un digne fils de S. Louis porte encore la noble couronne de France. C'est un prince solidement vertueux, également pénétré et de ses devoirs envers son Dieu et de ses devoirs envers son peuple. Mais son âme est trop belle pour le malheureux temps où il est appelé à régner. Comment pourrait-il soupçonner l'audace frénétique de ceux qui se sont faits ses ennemis ? Si son trône est encore debout depuis longtemps, hélas ! miné par la licence et l'irréligion, il ne ressemble que trop aux antiques débris qu'on croirait sains et entiers au milieu de ces cités jadis ensevelies sous la lave d'un volcan. Le moindre souffle va les faire tomber en poudre.

Déjà l'agitation commence ; les États généraux sont convoqués ; le premier acte de cette assemblée soi-disant nationale est un acte de révolte par lequel elle consacre le principe d'insurrection. Bientôt on attaque la religion par le schisme, en attendant que le moment soit venu de faire disparaître cette ombre sacrée qu'on lui conserve encore. On décrète la liberté pour tous ; mais cette liberté ainsi proclamée n'existe que pour le mal, elle est nulle pour le bien. Le roi est moins libre que personne ; il est livré à toutes les attaques des factieux et en butte à tous leurs outrages. Enfin le 10 août 1792, on achève de détruire ce qui semble rester encore de l'autorité royale. Le trône est renversé ; l'héritier d'une dynastie de huit siècles est jeté dans les fers. Quelques jours après, le 2 septembre, d'infâmes scélérats croient anéantir la religion en massacrant des évêques et des prêtres, prisonniers et sans défense ; mais ils n'ont réussi qu'à faire des martyrs. La révolution ne recule devant
% Source : PDF original, vue 007, côté droit, page manuscrite 13.
aucun crime. Une assemblée de régicides ose juger son souverain. Louis Seize, devant ces hommes iniques, parle avec la majesté d'un roi ; dans son immortel testament il tient le langage d'un saint, et le 21 janvier, victime d'expiation pour les crimes de son peuple, il meurt avec le courage d'un martyr.

Ceux qui n'ont pas craint de verser le sang de leur roi ne s'arrêteront pas là. Le règne de l'échafaud, si justement nommé le règne de la Terreur, commence et se personnifie dans le misérable, horriblement célèbre sous le nom à jamais détesté de Maximilien Robespierre.

\section*{Coup d'œil sur Lyon}
% Mention marginale de la page manuscrite 13.
L'impiété et l'anarchie révolutionnaire qui marchaient à sa suite avaient établi à Paris le siège de leur empire. Mais de là elles répandaient leurs poisons sur toutes les provinces. Notre ville n'en fut pas exempte ; aussi fut-elle violemment secouée par la tempête. Depuis longtemps l'éducation de la classe inférieure de la société y était très négligée. C'était un des tristes résultats de la destruction des Jésuites. Nous n'avions pas alors le bonheur de posséder les Frères de la doctrine chrétienne. Ces saints et modestes religieux auraient pu sous bien des rapports suppléer auprès de cette notable partie de la population lyonnaise à l'absence des enfans de S. Ignace. Avec cet ordre célèbre, dont les ennemis de la religion n'avaient que trop compris l'utilité, avaient aussi disparu du milieu de nous toutes les réunions pieuses connues sous le nom de Congrégations, au moyen desquelles des hommes zélés pour la gloire de Dieu savaient rallier autour d'eux,
% Source : PDF original, vue 008, côté gauche, page manuscrite 14.
pour les gagner à J.-C., les chrétiens de tous les états, et surtout cette multitude qui, dans toute la rigueur de l'expression, est obligée de gagner son pain à la sueur de son front. Une grande portion de cette classe (car la suite de ce récit prouvera qu'on pouvait compter beaucoup d'honorables exceptions), une grande partie donc n'avait conservé l'habitude des choses de la religion que comme une affaire de pure cérémonie ; et pour un grand nombre l'instruction sur les vérités du salut était à peu près nulle. Ceci explique comment il a été facile aux chefs de la conspiration formée contre toute espèce d'ordre public de les ameuter contre les prêtres, les nobles et les riches ; et de les lancer comme de véritables bêtes féroces sur la proie qu'ils offraient à leur avidité. Ainsi que nous l'avons déjà dit, il existait encore à Lyon, dans tous les états, une multitude de familles vraiment patriarcales, dans lesquelles la foi et l'amour de la religion étaient héréditaires. C'était surtout dans la classe moyenne qu'on en rencontrait davantage.

Les pères de famille avaient été pour la plupart élèves dans un des deux collèges que les Jésuites occupaient dans notre ville avant leur expulsion. Ils avaient transmis à leurs enfants les vieilles traditions qu'ils avaient puisées à une source si pure. De plus, par un usage général d'un heureux établissement, les jeunes gens, après avoir terminé leurs études de latin, allaient passer deux années au grand séminaire de St Irénée pour faire sous la direction des Sulpiciens leur cours de philosophie. Sous ces maîtres habiles, remplis de l'esprit de Dieu et connaissant assez
% Source : PDF original, vue 008, côté droit, page manuscrite 15.
les orages du siècle pour pouvoir diriger leurs élèves au moment où ils entraient dans le monde, ces jeunes gens voyaient s'affermir et se fortifier dans leurs cœurs les principes de la foi. Aussi cette portion éclairée de la jeunesse lyonnaise, bien loin de se laisser séduire, a montré une profonde aversion pour les maximes de l'incrédulité et un zèle fort ardent contre l'esprit révolutionnaire. Parmi elle on rencontrait des modèles de vertu et des cœurs animés de la piété la plus vive. Cette jeunesse qui dans l'âge mûr eût si fort honoré notre ville n'a été, hélas ! que trop décimée par le siège de Lyon et les temps affreux qui en ont été la suite. Néanmoins, parmi ceux qui ont pu échapper à la tourmente révolutionnaire, Dieu dans sa miséricorde en a destiné plusieurs à venir prendre place dans les rangs de la Congrégation, aussitôt qu'elle a été établie, et à l'aider de leur concours et de leurs lumières dans ces premiers et si faibles commencements. Nous avons été les témoins de leur piété et de leur zèle ; et nous pouvons signaler ici, comme une des causes de la prospérité que le Seigneur réservait à notre Congrégation, les bons sentimens qu'il avait placés et conservés dans les cœurs de ces hommes élèves du séminaire de St Irénée.

Reprenons maintenant le cours des événements. Dans les provinces et surtout à Lyon, toutes les choses de la Révolution, après s'être montrées à Paris, se reproduisaient par une affreuse imitation. Dès 1791, le schisme fut violemment établi dans notre ville. Les prêtres fidèles
% Source : PDF original, vue 009, côté gauche, page manuscrite 16.
à la voix de leur conscience, qui ne voulaient ni prêter serment à la Constitution civile du clergé, ni reconnaître l'évêque intrus Lamourette, furent arrachés à leurs troupeaux et forcés pour la plupart de s'enfuir.

Pour l'instruction de nos plus jeunes frères, dont quelques-uns peut-être ne connaissent qu'imparfaitement l'histoire de ces temps malheureux, je dois vous rappeler ici qu'on a appelé Constitution civile du clergé une loi par laquelle l'Assemblée dite nationale avait, sans le consentement du chef de l'Église, bouleversé toute la hiérarchie de l'Église de France, qu'elle prétendait rendre à peu près indépendante du St-Siège. Cette audacieuse entreprise avait été condamnée par le pape et par tous les évêques de France, moins quatre. On a désigné par le nom de réfractaires ou insermentés les prêtres catholiques qui, n'ayant pas voulu adhérer au schisme, ont refusé le serment exigé par le gouvernement. Ceux qui ont eu la foiblesse de faire ce serment étaient appelés prêtres constitutionnels ou assermentés. Enfin les évêques et les prêtres qui prenaient la place des pasteurs fidèles que la persécution arrachait à leur poste reçurent le nom d'évêques et de curés intrus.

La Révolution ne se borna pas seulement à chasser, par l'exécution de la loi de la déportation, les prêtres insermentés de la France. Outre le massacre du 1\up{er} septembre à Paris, dont nous avons déjà parlé, il y eut aussi d'autres prêtres qui obtinrent dans les provinces, et à Lyon en particulier dès le commencement de la
% Source : PDF original, vue 009, côté droit, page manuscrite 17.
persécution, la couronne du martyre. D'autres seulement emprisonnés pour la cause de la religion méritaient le titre de confesseurs.

\section*{1793 -- Terreur}
% Mention marginale de la page manuscrite 17.
On vit à cette époque, où les scènes de désordre semblaient se multiplier de plus en plus, à Lyon comme à Paris, une tourbe effrénée de fanatiques se jeter avec la rage des chiens enragés sur les personnes pieuses qui venaient assister aux offices dans les églises qui étaient encore desservies par des prêtres catholiques, et livre de foibles femmes à d'indignes outrages. On les vit plus tard organiser dans la ville le meurtre et le pillage. On préludait par ces horreurs à des crimes plus odieux encore, si toutefois cela est possible. On méditait d'étendre ce système d'assassinat dans lequel on était entré, de manière à y comprendre tous les citoyens à la fortune desquels on pouvait porter quelque envie. L'excès du mal amena la résistance. Le 29 mai 1793, l'anarchie fut vaincue dans les rues de Lyon ; mais malheureusement, le 31 mai, elle fut victorieuse à Paris. Les conséquences de sa victoire pesèrent cruellement sur nous. Les Jacobins (tel était le nom qu'on donnait alors à ces hommes de sang), les Jacobins rassemblent leurs satellites ; ils s'élancent en frémissant de rage sur la proie qui leur est échappée. Ils jurent de se venger et d'anéantir la cité rebelle. La jeunesse lyonnaise court aux armes ; par une défense héroïque elle arrête à nos portes pendant plus de deux mois ces bourreaux avides de se baigner dans notre sang. Il fallut enfin succomber !... Comment raconter maintenant à nos jeunes amis les scènes d'horreur
% Source : PDF original, vue 010, côté gauche, page manuscrite 18.
qui suivirent ! Ils ne connaissent ces temps de hideuse mémoire que par le souvenir qu'ont pu leur en conserver leurs pères ; ils vont me taxer d'exagération, et cependant ma plume se refuse à tout dire. Comment en effet reproduire à vos yeux un tableau aussi lamentable ? Les échafauds non seulement dressés, mais fixes solidement sur les places publiques !... Le sang versé chaque jour par torrent. La balle meurtrière suppléant à la lenteur du fer homicide, trop peu expéditif au gré de nos affreux tyrans ! Le canon pointé avec une joie féroce contre des troupes nombreuses de victimes, venant accélérer encore l'effroyable rapidité du carnage !... Vainement quelqu'un du grand nombre conservera l'espérance d'échapper à la faux de la mort qui semble, par mégarde, l'avoir oublié ; des bourreaux organisés en légions, pour cet odieux service, se précipitent sur cet horrible champ de bataille et achèvent à coups de sabre les victimes qui palpitent encore ou celles que la mitraille avait épargnées. Toute marque de deuil, tout signe de piété, la moindre sympathie à l'égard des malheureux immolés à ce qu'on appelle la vengeance nationale, est un crime irrémissible jugé digne de mort. Voilà, Messieurs, une bien foible esquisse de ces temps de sanglante mémoire qu'on a appelés la Terreur.

Mais au milieu de toutes ces horreurs, que faisait la religion ? Ce qu'elle faisait, Messieurs ? Nous la retrouverons ici comme on la retrouve dans tous les temps et partout, toujours noble et grande. Appuyée sur la croix de J.-C., elle se tenait debout,
% Source : PDF original, vue 010, côté droit, page manuscrite 19.
bravant la tempête. Elle consolait les fidèles en offrant dans le secret l'adorable victime d'expiation, en présence de chrétiens fervents que la grâce d'en haut remplissait de force et de courage, car ils savaient que leur assistance à la messe suffirait pour les conduire au supplice. Pour ouvrir le ciel par une dernière absolution à ceux sur la tête desquels va se lever le fer des bourreaux, elle envoyait ses ministres dans les prisons et jusque sur les marches de l'échafaud, que bientôt peut-être ils teindront de leur sang. Elle conduisait au ciel, par le chemin qu'a tenu le Sauveur du monde et qu'ont suivi de nombreux athlètes dans les premiers siècles, des légions de nouveaux martyrs.

Parmi eux et en première ligne paraissent ces prêtres du Seigneur, ministres de charité qui venaient avec joie terminer sur l'échafaud une vie entièrement consacrée à la défense de la foi, en rendant par leur mort un nouveau témoignage à la vérité de ses oracles. Parmi eux on retrouvait de nouvelles Félicité, de nouvelles Perpétue, de nouvelles Blandine, le cœur tout pénétré de l'amour de leur divin maître, marchant à la mort comme à la victoire et montrant par la grandeur de leur héroïsme aux spectateurs étonnés comment la grâce de N. S. J.-C. sait faire rencontrer la force jusque dans la foiblesse. Parmi eux apparaît encore de nouveaux Éléazar, martyrs de la vérité qui, plutôt que de souiller leur bouche par un mensonge, préféraient glorifier par leur mort le Dieu de toute vérité. Parmi eux enfin, Messieurs, se sont rencontrés des parens,
% Source : PDF original, vue 011, côté gauche, page manuscrite 20.
des amis de plusieurs d'entre nous. Si je traçais ici des noms chers à plusieurs de nos frères, je ferais palpiter bien des cœurs. C'est ainsi que dans la première page de ce récit, longtemps avant que le moment soit arrivé où la Congrégation devait être fondée, je vous raconte déjà votre propre histoire. Ce sang si pur versé pour la cause de la foi, c'était le sang de nos concitoyens, c'était notre propre sang, celui de plusieurs membres de nos familles. Ces martyrs de la fin du 18\up{e} siècle sont allés rejoindre au pied du trône de l'Éternel nos martyrs du second siècle. Tous ensemble ont fait en notre faveur une sainte violence au ciel. Le sang des martyrs est toujours une semence féconde ; il est redescendu sur nous comme une rosée abondante, et peut-être est-il devenu le germe précieux qui a produit la Congrégation. Soyons-en bien convaincus, c'est à l'intervention de nos martyrs que nous avons dû toutes ces grâces de prédilection dont nous avons reçu les prémices ; ce sont eux qui nous ont obtenu la protection de Marie, la reine des martyrs.

C'est ainsi que l'œuvre de Dieu va toujours s'accomplissant au milieu des calamités qui paraissent les plus interminables. Cette terreur qui couvrait notre patrie d'un voile de deuil horriblement ensanglanté avait pour but, dans les décrets du Très-Haut, tout en punissant la France, de peupler le ciel de nouveaux saints et de faire pousser à l'arbre de la foi, en l'arrosant de leur sang, de profondes racines au milieu de nous.

% Source : PDF original, vue 011, côté droit, page manuscrite 21, première partie.
% Le chapitre second commence au milieu de cette même page.
\section*{1794}
Cette funeste période de notre Révolution se prolongea jusqu'au mois de juillet 1794. À cette époque, l'infâme Robespierre, ce monstre abominable qui régnait en souverain par lui ou par ses délégués sur notre pays désolé, fut prévenu par plusieurs de ses complices, au moment où il avait formé le projet de les perdre. Abandonné de tous les siens, il est arrêté et reçoit en cherchant à se défendre une blessure qui le couvre de sang ; et est ainsi porté tout sanglant sur l'échafaud où l'accompagnent les malédictions et les injures d'une population révoltée par tant d'horreurs. Cette journée fut saluée par la France comme le signal de la délivrance. On l'appela la journée du neuf Thermidor.
